\providecommand{\itescastudentname}{Martin Jonathan de la Cruz}
\providecommand{\itescafaculty}{Maestría en Gestión Administrativa}
\providecommand{\itescacoursename}{Seminario I}
\providecommand{\itescateachingfigure}{Dra. Carla Olimpya Zapuche Moreno}
\providecommand{\itescalocation}{Cajeme, Sonora}
\providecommand{\itescaofficialasset}{ITESCA/assets-itesca/web/logo-itesca-contorno.png}
\providecommand{\itescabibliography}{ITESCA/maestria-en-gestion-administrativa/seminario-i-mga/seminario-i}

\documentclass[12pt,spanish,letterpaper,oneside]{article}
\usepackage[utf8]{inputenc}
\usepackage[T1]{fontenc}
\usepackage[spanish,mexico]{babel}
\usepackage[letterpaper,margin=2.54cm]{geometry}
\usepackage{helvet}
\renewcommand{\familydefault}{\sfdefault}
\usepackage{setspace}
\usepackage{graphicx}
\usepackage{csquotes}
\usepackage{natbib}
\usepackage{bibentry}
\usepackage{url}
\usepackage[hidelinks]{hyperref}
\usepackage{ragged2e}
\usepackage{enumitem}
\usepackage{fancyhdr}

\setlength{\parindent}{1.27cm}
\setlength{\parskip}{0pt}
\doublespacing
\urlstyle{same}
\setlist{nosep}
\pagestyle{fancy}
\fancyhf{}
\fancyhead[R]{\thepage}
\renewcommand{\headrulewidth}{0pt}
\setlength{\headheight}{15pt}

\newenvironment{apablockquote}{%
  \begin{list}{}{\setlength{\leftmargin}{1.27cm}\setlength{\rightmargin}{0cm}}\item\relax
}{\end{list}}

\newenvironment{apareferences}{%
  \begin{list}{}{%
    \setlength{\leftmargin}{1.27cm}%
    \setlength{\itemindent}{-1.27cm}%
    \setlength{\itemsep}{0pt}%
    \setlength{\parsep}{0pt}%
  }%
}{\end{list}}

\begin{document}

\begin{titlepage}
\thispagestyle{empty}
\centering
\vspace*{1.2cm}
\includegraphics[width=4.5cm]{\itescaofficialasset}

\vspace{1.8cm}
{\large Instituto Tecnológico Superior de Cajeme\par}
\vspace{2.2cm}
{\bfseries\Large Mi primera práctica en formato APA\par}
\vspace{0.5cm}
{\large Actividad 2\par}
\vfill
\begin{flushleft}
\setlength{\parindent}{0pt}
\textbf{Nombre completo:} \itescastudentname\\
\textbf{Programa de estudios:} \itescafaculty\\
\textbf{Asignatura:} \itescacoursename\\
\textbf{Docente:} \itescateachingfigure\\
\textbf{Fecha:} 29 de agosto de 2026
\end{flushleft}
\vfill
{\itescalocation\par}
\end{titlepage}

\setcounter{page}{2}

\section{Introducción}

\textbf{Objetivo:} aplicar las reglas de citación y referencia de APA 7 en un argumento coherente sobre decisiones administrativas basadas en evidencia. El problema central es que una decisión puede parecer razonable y, aun así, carecer de datos verificables o atribuir como propia una idea ajena. Por ello, el criterio de evaluación exige distinguir una paráfrasis, una cita breve y una cita en bloque, y vincular cada una con una referencia recuperable.

\section{Evidencia y decisiones en la gestión administrativa}

\subsection{Práctica de paráfrasis y cita textual corta}

La gestión administrativa basada en evidencia permite sustituir decisiones apoyadas únicamente en la intuición por juicios que integran datos de la organización, conocimiento científico y experiencia profesional. \citet{rousseau2006Evidence} sostiene que este enfoque convierte los mejores hechos disponibles en una base para las prácticas directivas; en el mismo sentido, \citet{hernandezSampieri2014Metodologia} explican que la investigación sistemática articula procesos críticos y empíricos para estudiar problemas, por lo que sus resultados pueden orientar diagnósticos y decisiones sin confundir una opinión con una conclusión sustentada. Esta perspectiva es especialmente útil en mercadotecnia y ventas: antes de modificar una campaña, un proceso o una meta, la persona responsable debe formular una pregunta, comparar indicadores, valorar la calidad de las fuentes y documentar el criterio elegido. La complejidad cotidiana no desaparece por seguir una secuencia formal; como advierten \citet{brightCortes2019Principles}, el mundo real de quienes administran es un \textquote{flujo desordenado y frenético de actividad continua} (cap. 1, párr. 4, traducción propia). La cita corta resume una dificultad concreta, mientras que la paráfrasis anterior integra ideas ajenas con una redacción propia y conserva la atribución correspondiente.

\subsection{Práctica de cita textual larga}

La necesidad de interpretar la evidencia dentro de su contexto también se aprecia en la siguiente descripción del trabajo directivo:

\begin{apablockquote}
La mayoría de los libros de administración dirían, como lo hace este, que los gerentes dedican su tiempo a planear, organizar, dotar de personal, dirigir, coordinar, informar y controlar. Estas actividades, como encontró Hannaway en su estudio de los gerentes en el trabajo, 
\textquote{no describen, de hecho, lo que hacen los gerentes}. En el mejor de los casos parecen describir objetivos vagos que los gerentes intentan alcanzar continuamente. (Bright \& Cortes, 2019, cap. 1, párr. 4, traducción propia)
\end{apablockquote}

El pasaje supera 40 palabras y, por ello, se presenta como cita en bloque, sin comillas externas, con sangría y localizador. Su contenido muestra que una lista de funciones no basta para explicar una decisión administrativa: es necesario reconstruir el problema, las restricciones y los efectos esperados. Un antecedente académico aplicado examinó precisamente la toma de decisiones en la gestión administrativa de una empresa \citep{zambrano2018Decisiones}, mientras que \citet{sinek2009Leaders} propuso, en una conferencia, comunicar primero el propósito para movilizar la acción. Ambos recursos cumplen funciones distintas: una tesis aporta un estudio documentado y el video ofrece una exposición profesional que debe evaluarse de acuerdo con su naturaleza. En consecuencia, una práctica responsable combina fuentes pertinentes, explicita qué afirmación respalda cada una y aplica las reglas de citación: la American Psychological Association indica que las citas menores de 40 palabras se incorporan entre comillas y que las de 40 palabras o más se disponen en bloque \citep{apa2022Quotations}.

\section{Conclusiones}

Considero que el uso correcto de APA no es un adorno tipográfico, sino un mecanismo de integridad académica. Mi postura se sustenta en que una decisión administrativa solo puede someterse a crítica cuando el lector distingue qué parte corresponde al análisis propio y cuál procede de una fuente verificable. En consecuencia, citar con precisión fortalece la argumentación, reduce atribuciones ambiguas y establece una práctica transferible a la futura tesis y a los informes profesionales.\footnote{Se utilizó inteligencia artificial como apoyo para organizar ideas y revisar la redacción; la selección de fuentes, su verificación y el análisis final fueron realizados y supervisados por el autor.}

\clearpage
\section{Referencias}

\bibliographystyle{apalike}
\bibliography{\itescabibliography}

\end{document}
